% ---------------------------------------------------------------
% Section 2 draft: Geometry and model representation
% Replaces the current "Drafting and Model Representation" section.
% Numbers below are for the baseline (all d = 286 mm, R_f = 50 mm).
% ---------------------------------------------------------------
\section{Geometry and Model Representation}

\subsection{Parametric construction}
The frame profile is generated by a single FreeCAD~1.1 Python script
(\texttt{rebuild\_truss.py}, Appendix~A) built on the OpenCASCADE kernel.
The script's only inputs are the node coordinates, one width
$d_e\in\{89,140,184,235,286\}$~mm per member, the internal-corner radius
$R_f$, the support extension $e=50$~mm and the thickness $b=76$~mm.
Every candidate design is therefore defined by a short parameter block
rather than a hand-drawn sketch. All lengths are in mm. The origin is
at $B_0$, $x$ runs along the span and $y$ is vertical; the eaves are at
$y=0$ and the ridge is at $R=(4000,\,1333.3)$~mm (4:12 pitch,
$\theta=18.43^\circ$). The geometry is built in five steps
(Fig.~\ref{fig:geom}):

\begin{enumerate}
  \item \textbf{Member strips.} Each member $e$ is a rectangle of width
  $d_e$ centered on its node-to-node centerline. The two lower-chord
  strips are lengthened by $e=50$~mm past $B_0$ and $B_2$ to form the
  support segments.
  \item \textbf{Joint discs.} At every node a circular disc is added
  whose diameter equals the widest member meeting there. The disc covers
  the square strip ends, so no corner of a strip projects past the outline
  of the joint (for example, at the ridge the disc of radius 143~mm is
  tangent to both upper-chord faces).
  \item \textbf{Boolean union and trim.} All strips and discs are fused
  into one face, so overlapping joint material is counted once. The face
  is then clipped by a flat line along the lower face of the lower chord
  ($y=-d/2$) and by vertical lines at $x=-50$ and $x=8050$~mm. This leaves
  two flat, constant-width support segments
  $x\in[-50,50]$ and $x\in[7950,8050]$~mm on the lower surface.
  \item \textbf{Internal-corner rounding.} All re-entrant corners are
  rounded with $R_f=50$~mm ($\geq20$~mm required) by morphological
  closing: the outer boundary is offset by $+R_f$ and then $-R_f$, and
  each of the four openings by $-R_f$ and then $+R_f$, using arc joins.
  This rounds only concave material corners and leaves convex corners
  sharp. The rounding geometry is part of the face and so is included in
  the volume.
  \item \textbf{Export.} The script checks that the face is a single valid
  region with the expected number of openings. It reports
  $V=b\,|\Omega|$, then writes the 2D face as a BREP for meshing in gmsh
  and an extruded solid as an \texttt{.FCStd} file for visualization.
\end{enumerate}

\subsection{Joint representation}
Each joint is a continuous solid region (disc plus filleted overlap),
which represents an ideal bonded connection. Members are connected only
because they share material in this single face, which then gets one
conforming mesh; no coincident-node ties are used. Fasteners, plates,
bearing and slip are not modeled. The same joint construction
(widest-member disc, $R_f=50$~mm) is used for every design that is
compared.

\subsection{Baseline geometry check}
For the baseline (nine members, all $d=286$~mm) the profile area is
$|\Omega|=5.622\times10^{6}$~mm$^2$. This gives $V=0.427$~m$^3$ and
$\rho V=214$~kg per frame, or $21V=8.98$~m$^3$ and 4490~kg for the
whole roof, compared with the centerline estimate
$\sum b d_e \ell_e=0.478$~m$^3$. The difference is the overlap at
joints, which is counted only once in the union. The fillets add
$0.061\times10^{6}$~mm$^2$ (1.1\%). Rebuilding the profile independently
from the same parameters with Shapely gives the same area.

The upper and side faces stay within 143~mm of the centerline triangle,
with the ridge disc reaching $y=1476.3$~mm. The only material beyond
143~mm is at the four corners of the 50~mm support extensions
($(-50,\pm143)$ and its mirror, 151.5~mm from $B_0$ and $B_2$). We treat
these extensions as required by the support specification rather than
as member width.

\subsection{Boundary identification}
The support boundaries are the edges at $y=-d/2$ with
$|x-x_{B}|\leq50$~mm. The loaded upper boundary $\Gamma_q$ consists of
the exterior edges with an upward outward normal. Its horizontal
projection, computed from the mesh, is 8100~mm, which includes the two
50~mm extensions. [\emph{State your choice: either restrict $q$ to
$0\leq x\leq8000$~mm with an indicator in the UFL form so that the total
load matches the 8.0~m projected area used in the truss, or keep
8100~mm and apply the same total in the truss.}]
